\chapter{Análisis del problema}
 
El problema abordado en este trabajo consiste en diseñar e implementar una herramienta web capaz de visualizar y explorar grafos cerebrales (connectomes) de forma interactiva en tres dimensiones. Aunque existen soluciones de escritorio en el ámbito de la neuroimagen, muchas requieren entornos complejos de instalación o no están orientadas a una experiencia de uso inmediata desde navegador.

En este contexto, el reto principal no es solo representar nodos y aristas, sino hacerlo de forma comprensible para el usuario final, manteniendo un equilibrio entre fidelidad de los datos, rendimiento gráfico y simplicidad de interacción.

\section{Contexto y motivación}

Un connectome se modela habitualmente como un grafo ponderado donde:

\begin{itemize}
	\item cada nodo representa una región cerebral,
	\item cada arista representa una conexión estructural o funcional,
	\item el peso de la arista expresa intensidad o relevancia de la conexión.
\end{itemize}

La combinación de alta densidad de conexiones y representación espacial 3D introduce dificultades claras para el análisis visual: oclusión, sobrecarga perceptiva y pérdida de contexto. Por ello, resulta necesario proporcionar mecanismos de filtrado, resaltado y navegación que permitan reducir la complejidad sin eliminar información relevante.

\section{Definición del problema}

El sistema debe transformar dos fuentes de datos principales (fichero de nodos y matriz de adyacencia) en una representación interactiva navegable en tiempo real desde navegador web. Esta transformación implica varias subtareas:

\begin{enumerate}
	\item Parseo robusto de ficheros heterogéneos.
	\item Validación de consistencia entre nodos y aristas.
	\item Exposición de datos mediante una API REST estable.
	\item Renderizado 3D eficiente en cliente.
	\item Interacciones de usuario orientadas al análisis exploratorio.
\end{enumerate}

Desde el punto de vista técnico, el problema se puede resumir como una tubería de procesamiento y visualización donde cualquier error en la etapa de datos se propaga al renderizado final.

\section{Requisitos de análisis}

Para que la solución sea útil en un entorno académico y reproducible, se establecen los siguientes requisitos de análisis:

\begin{itemize}
	\item \textbf{Exactitud de datos:} los enlaces visualizados deben corresponder con la estructura del grafo de entrada.
	\item \textbf{Interpretabilidad:} el usuario debe poder identificar regiones, grupos y relaciones sin ambigüedad.
	\item \textbf{Interactividad:} la navegación por cámara, el hover y la selección deben responder con baja latencia.
	\item \textbf{Escalabilidad moderada:} la aplicación debe soportar datasets de tamaño medio sin degradación crítica.
	\item \textbf{Despliegue sencillo:} el sistema debe levantarse con un flujo reproducible basado en contenedores.
\end{itemize}

\section{Restricciones y supuestos}

El análisis también contempla limitaciones realistas del proyecto:

\begin{itemize}
	\item El frontend se ejecuta sobre WebGL, por lo que el rendimiento depende de la capacidad gráfica del equipo cliente.
	\item La versión inicial prioriza un dataset de referencia (AAL90), por lo que la generalización a grafos mucho mayores requerirá técnicas adicionales (LOD, clustering dinámico o simplificación de geometría).
	\item El backend se orienta a servir datos preprocesados y métricas básicas, no a ejecutar pipelines de neuroimagen pesados.
	\item El alcance de autenticación y persistencia avanzada queda fuera del núcleo MVP y se considera evolución posterior.
\end{itemize}

\section{Riesgos técnicos identificados}

Durante el análisis se identifican los siguientes riesgos:

\begin{itemize}
	\item \textbf{Sobrecarga visual:} demasiadas aristas visibles simultáneamente reducen la utilidad analítica.
	\item \textbf{Inconsistencia de formatos:} pequeñas variaciones en los ficheros de entrada pueden romper el parseo.
	\item \textbf{Cuellos de botella en renderizado:} número elevado de draw calls o geometrías no optimizadas.
	\item \textbf{Dependencia del entorno local:} conflictos de puertos o diferencias de sistema operativo en desarrollo.
\end{itemize}

Como mitigación, se adopta una arquitectura desacoplada (API + cliente 3D) y una estrategia de despliegue con Docker Compose, reduciendo variabilidad del entorno.

\section{Criterios de éxito}

Se considera que el problema está correctamente resuelto en su alcance actual cuando:

\begin{itemize}
	\item el sistema carga y muestra el connectome completo sin errores,
	\item el usuario puede explorar la escena 3D de forma fluida,
	\item existe trazabilidad entre datos de entrada y visualización final,
	\item el despliegue es reproducible mediante una única secuencia de comandos.
\end{itemize}

Este análisis fundamenta las decisiones de diseño presentadas en los capítulos siguientes y justifica la elección de tecnologías empleadas en la implementación.
 
